% Chapter 1 - Introduction

\section{Introduction}\label{sec:introduction}

The starting point for this thesis was a simple question: can you make a
hybrid game and editing application driven by hand motion using just a webcam and
a regular \acrshort{cpu}, with no extra hardware at all? No \acrshort{vr} headset,
no depth camera, no special gloves --- just the camera that is already built into a laptop.
This idea came from wanting to make \acrshort{hci} as accessible as possible.
Expensive peripherals limit who can use a system, and the setup time alone is
often a reason people do not bother. The human hand is the most natural
controller there is, and ideally it should work out of the box.

This is now within reach thanks to frameworks like Google MediaPipe
\cite{zhang2020mediapipehands}, which can detect both hands in real time
from a standard colour camera, outputting 21 2.5D landmarks per hand ---
comprising x, y, and relative depth coordinates --- at well under 20 milliseconds
per frame on modern mobile devices such as the Pixel 3, Samsung S20, and iPhone 11.
The problem, however, is that these raw landmarks are not directly
useful for driving a skeleton or analysing motion. Their scale changes with
how close the hand is to the camera, depth is estimated rather than measured,
and there is no information about joint hierarchy or anatomically valid
ranges of motion --- it is just a flat list of points in space.

\subsection{The Problem This Thesis Addresses}\label{subsec:gap}

To actually drive a hand avatar or store motion data in a meaningful way,
you need a \emph{canonical} representation: something that describes the
hand as a set of joint rotations along a fixed kinematic tree, independent
of hand size or camera distance. The MANO parametric hand model
\cite{MANO:SIGGRAPHASIA:2017} is exactly this. It represents hand pose
as 15 joint rotations plus a global orientation, in a predefined parent-child
kinematic tree whose parameters (blend shapes, skinning weights, joint regressor)
were learned from real hand scans. A pose in MANO space can drive a 3D mesh,
be logged, compared across sessions, or passed to a classifier without
any additional preprocessing.

The core challenge of this thesis is therefore to bridge the gap between
MediaPipe's raw output and MANO's canonical representation, in real time,
on a CPU, and with low enough latency to feel responsive in an interactive
application.

\subsection{Contributions}\label{subsec:contributions}

The work presented here covers the full pipeline from webcam input to
a running Unity application. Specifically:

A hand tracking and 3D coordinate module (\texttt{hand\_engine},
\texttt{hand\_kinematics\_3D}) retrieves metric landmarks from MediaPipe
and filters them using the 1\euro{} filter \cite{casiez2012oneeuro} to remove
jitter while keeping latency low.

An Adaptive \acrshort{ik} solver (\texttt{MANOkinematics}) maps the 21
MediaPipe landmarks onto MANO's 16-joint skeleton using the AIK algorithm
\cite{zhou2020monocular} \cite{meng2021aik}, producing axis-angle joint rotations that are invariant
to hand scale and camera distance.

A UDP streaming layer (\texttt{mano\_unity\_streamer},
\texttt{ManoLiveReceiver.cs}) sends the solved pose and detected gesture
for both hands to Unity as a JSON datagram, achieving a measured mean
software latency of \SI{30.68}{\ms} over more than 20\,000 frames.

A Unity prototype (\texttt{my\_hands\_game}) with four scenes that each
use the skeletal stream differently: a physics-based grapple interaction,
an editable city environment, a spring-physics face deformer, and a 2D/3D
painting tool with point-cloud export.
